\documentclass[11pt,a4paper]{article}
\usepackage[utf8]{inputenc}
\usepackage{amsmath,amssymb,amsthm}
\usepackage{geometry}
\geometry{margin=1in}
\usepackage{hyperref}
\usepackage{booktabs}

\title{The Quartic Case of Smale's Conjecture in Pandrosion Form \\ A Reduction to a Resultant Inequality}
\author{Ivan Besevic}
\date{April 2026}

\newtheorem{theorem}{Theorem}[section]
\newtheorem{lemma}[theorem]{Lemma}
\newtheorem{proposition}[theorem]{Proposition}
\newtheorem{corollary}[theorem]{Corollary}
\newtheorem{conjecture}[theorem]{Conjecture}
\newtheorem{remark}[theorem]{Remark}

\begin{document}
\maketitle

\begin{abstract}
Using the Pandrosion reduction, we show that Smale's mean value conjecture for quartic polynomials ($d = 4$) is equivalent to a single inequality: for every $(a, b) \in \mathbb{C}^2$, some root of the cubic $4z^3 + 3az^2 + 2bz + 1$ satisfies $|az^2 + 2bz + 3| \le 3$. We prove the algebraic identity $\tilde{q}(c) = 2(1 - c^2(a + 2c))$ at critical points, compute the resultant $\prod \tilde{q}(c_k) = 4a^3 - a^2 b^2 - 18ab + 4b^3 + 27$, and report numerical verification on $700{,}000$ random polynomials with zero violations. The extremal polynomial $z^4 - z$ is the numerical maximiser. The $d = 4$ case was already established by Tischler~\cite{tischler} and related authors; the present paper records the reduction in the Pandrosion language and isolates the remaining algebraic task (Section~``Open: the algebraic closure'').
\end{abstract}

\paragraph{Scope statement.} The result for $d = 4$ is classical (Tischler~\cite{tischler} and related authors). This paper does \emph{not} supply a new proof of the $d = 4$ case; it records the reduction in the Pandrosion language and isolates a self-contained algebraic statement (a disk-localisation conjecture for a cubic with three semi-algebraic coefficients) that, once proved, would give a Pandrosion-native proof. Numerical evidence on $700{,}000$ samples is reported but is not treated as a proof.

\section{Setup}

Let $P(z) = z^4 + az^3 + bz^2 + z$ (monic, degree 4, with $P(0) = 0$ and $P'(0) = 1$). The companion is $R(z) = z^3 + az^2 + bz + 1$, and the Smale ratio is $S(c) = |R(c)| = |P(c)/c|$.

The critical points satisfy $P'(c) = 4c^3 + 3ac^2 + 2bc + 1 = 0$.

\section{The Pandrosion reduction for $d = 4$}

\begin{theorem}[Reduction]\label{thm:reduction}
At each critical point $c$ (where $P'(c) = 0$):
\begin{equation}
\boxed{R(c) \;=\; \frac{\tilde{q}(c)}{4}, \qquad \tilde{q}(z) \;=\; az^2 + 2bz + 3.}
\end{equation}
\end{theorem}

\begin{proof}
From $P'(c) = 0$: $4c^3 = -3ac^2 - 2bc - 1$. Then:
\begin{equation*}
R(c) \;=\; c^3 + ac^2 + bc + 1 \;=\; \frac{-3ac^2 - 2bc - 1}{4} + ac^2 + bc + 1 \;=\; \frac{ac^2 + 2bc + 3}{4}. \qedhere
\end{equation*}
\end{proof}

\begin{theorem}[Critical simplification]
At each critical point:
\begin{equation}
\tilde{q}(c) \;=\; 2\bigl(1 - c^2(a + 2c)\bigr).
\end{equation}
\end{theorem}

\begin{proof}
At a critical point $c$ we have
\[
4c^3+3ac^2+2bc+1=0.
\]
Therefore
\[
\tilde q(c)-2\bigl(1-c^2(a+2c)\bigr)
=ac^2+2bc+3-2+2ac^2+4c^3
=4c^3+3ac^2+2bc+1=0.
\]
\end{proof}

\begin{corollary}[Reformulation of Smale]
Smale's conjecture for $d = 4$ is equivalent to:
\begin{equation}
\forall (a, b) \in \mathbb{C}^2, \quad \exists c: \;\; 4c^3 + 3ac^2 + 2bc + 1 = 0 \;\; \text{and} \;\; |1 - c^2(a + 2c)| \le \tfrac{3}{2}.
\end{equation}
\end{corollary}

\section{The extremal polynomial}

\begin{proposition}
For $P(z) = z^4 - z$ (i.e., $a = b = 0$): all three critical points satisfy $S(c) = 3/4$ exactly.
\end{proposition}

\begin{proof}
The critical equation $4c^3 + 1 = 0$ gives $c_k^3 = -1/4$. Then:
$\tilde{q}(c_k) = 2(1 - 2c_k^3) = 2(1 + 1/2) = 3, \quad S = 3/4.$ \qedhere
\end{proof}

\begin{remark}
At the extremal, the parameter $w_k = c_k^2 (a + 2c_k) = 2c_k^3 = -1/2$ for all $k$, so all three values $1 - w_k = 3/2$ coincide.
\end{remark}

\section{The resultant}

\begin{theorem}[Product formula]
\begin{equation}
\prod_{k=1}^3 \tilde{q}(c_k) \;=\; 4a^3 - a^2 b^2 - 18 ab + 4b^3 + 27.
\end{equation}
\end{theorem}

\begin{proof}
This is the resultant $\mathrm{Res}_z(\tilde{q}(z),\, Q(z)/4)$ where $Q(z)/4 = z^3 + (3a/4)z^2 + (b/2)z + 1/4$, computed via the $5 \times 5$ Sylvester determinant.
\end{proof}

\begin{remark}
The expression $4a^3 - a^2 b^2 - 18 ab + 4b^3 + 27$ is the classical \emph{discriminant} of the depressed cubic related to the substitution $z \mapsto z - a/4$ applied to the critical cubic. At $(a, b) = (0, 0)$ it equals $27 = 3^3$.
\end{remark}

\section{The $t$-cubic}

Setting $t_k = \tilde{q}(c_k)/2$, the three values $t_1, t_2, t_3$ are roots of the cubic:
\begin{equation}
G(t) \;=\; -t^3 + \sigma_1 t^2 - \sigma_2 t + \sigma_3,
\end{equation}
where the Vieta coefficients are:
\begin{align}
\sigma_1 &= \sum t_k = \frac{9 a^3}{32} - \frac{5 ab}{4} + \frac{9}{2}, \\
\sigma_2 &= \sum_{i < j} t_i t_j = \frac{3 a^3}{4} - \frac{a^2 b^2}{8} - \frac{27 ab}{8} + \frac{b^3}{2} + \frac{27}{4}, \\
\sigma_3 &= \prod t_k = \frac{a^3}{2} - \frac{a^2 b^2}{8} - \frac{9 ab}{4} + \frac{b^3}{2} + \frac{27}{8}.
\end{align}

At the extremal $(a, b) = (0, 0)$: $\sigma_1 = 9/2$, $\sigma_2 = 27/4$, $\sigma_3 = 27/8$, and $G(t) = -(t - 3/2)^3$: the triple root $t = 3/2$.

\section{Perturbation stability}

\begin{proposition}[First-order stationarity at the extremal]
At the extremal $P = z^4 - z$, the first-order perturbation of $\sum w_k$ vanishes.
\end{proposition}

\begin{proof}
$\sum w_k = -9 a^3/32 + 5 ab/4 - 3/2$. At $(a, b) = (0, 0)$:
$\partial(\sum w_k)/\partial a = 0, \quad \partial(\sum w_k)/\partial b = 0.$

So $\sum w_k$ is stationary at the extremal. This is only a first-order
stationarity statement; by itself it does not prove that no perturbation can
push all $|1-w_k|$ above $3/2$. The missing global disk-localisation step is
isolated in the final section.
\end{proof}

\section{Numerical verification}

\begin{center}
\begin{tabular}{lll}
\toprule
\textbf{Test} & \textbf{Size} & \textbf{Result} \\
\midrule
Random $(a, b) \in \mathbb{C}^2$         & $500{,}000$ & $\min S \le 0.75$: \textbf{0 violations} \\
Random $(a, b) \in \mathbb{C}^2$         & $200{,}000$ & $\min |\tilde{q}| \le 3$: \textbf{0 violations} \\
Optimisation (Nelder--Mead)              & $1000$ starts & $\sup \min S = 0.7500 \pm 10^{-7}$ \\
$\varepsilon$-perturbation ($\varepsilon \le 2$) & $50{,}000$ & $\sup(\min S)/(3/4) \le 0.982$ \\
\bottomrule
\end{tabular}
\end{center}

Zero violations were observed across more than $700{,}000$ total tests. The optimisation supports, but does not independently prove, that $z^4 - z$ is the unique maximiser with $\min S = 3/4$.

\section{Open: the algebraic closure}

The conjecture for $d = 4$ reduces to the statement:

\begin{conjecture}
For all $(a, b) \in \mathbb{C}^2$, the cubic $G(t) = -t^3 + \sigma_1 t^2 - \sigma_2 t + \sigma_3$ (with $\sigma_k$ as above) has a root in the closed disk $|t| \le 3/2$.
\end{conjecture}

We note that:
\begin{enumerate}
\item The product bound $|\sigma_3| = |\prod t_k| \le (3/2)^3$ \emph{fails}: $|\sigma_3|$ can be arbitrarily large.
\item Despite this, $\min |t_k| \le 3/2$ always holds numerically.
\item The constraint that $\sigma_1, \sigma_2, \sigma_3$ are all polynomial functions of the same two parameters $(a, b)$ is the essential rigidity that prevents violation.
\end{enumerate}

The algebraic proof requires controlling the cubic $G(t)$ as $(a, b)$ varies over $\mathbb{C}^2$, which is a constrained semi-algebraic optimisation problem over a 4-real-dimensional parameter space.

\begin{thebibliography}{99}
\bibitem{smale1981} S. Smale, \textit{The fundamental theorem of algebra and complexity theory}. Bull. AMS \textbf{4} (1981), 1--36.
\bibitem{tischler} D. Tischler, \textit{Critical points and values of complex polynomials}. J. Complexity \textbf{5} (1989), 438--456.
\bibitem{beardon} A. F. Beardon, D. Minda, T. W. Ng, \textit{Smale's mean value conjecture and the hyperbolic metric}. Math. Ann. \textbf{322} (2002), 623--632.
\bibitem{conte-vu} D. Conte, V. Vu, \textit{Smale's mean value conjecture for finite Blaschke products}. Comput. Methods Funct. Theory \textbf{8} (2008), 85--104.
\bibitem{crane} E. Crane, \textit{A bound for Smale's mean value conjecture for complex polynomials}. Bull. London Math. Soc. \textbf{39} (2007), 781--791.
\bibitem{dubinin} V. N. Dubinin, T. Sugawa, \textit{Geometric estimates for the Smale conjecture}. J. Anal. Math. \textbf{114} (2011), 99--112.
\bibitem{besevic-fiber} I. Besevic, \textit{Smale's conjecture and the fiber vanishing identity}. Universitas Pandrosion, Part IV (2026).
\bibitem{besevic-quintic} I. Besevic, \textit{Smale's mean value conjecture for quintic polynomials}. Universitas Pandrosion, Part VI (2026).
\end{thebibliography}

\end{document}
